\subsection{Funciones disponibles y no disponibles sin conexión}
\label{sec:funciones-offline}

Conforme a RT-03.10, la captura del turno de 14 horas en terreno y la operación local de 24 horas en los centros de distribución conservan sus eventos hasta recibir confirmación. RT-03.13 exige nombrar también lo que no funciona sin red y el procedimiento que lo suple. El Anexo 4.1-J distingue ambos casos; una operación pendiente nunca equivale a aprobación de un tercero.


La operación local conserva la captura; las autorizaciones externas esperan conexión o procedimiento aprobado. AL-OFF-01, en el Anexo 4.1-M, exige un corte de 24 horas con relevos a las 8 y 16 horas, verificación de permisos, reinicio de componentes y reconciliación. La prueba incluye recepción, preparación y despacho con documentación tributaria disponible; también ensaya el despacho sin guía, cuya brecha se resuelve mediante AL-DTE-01. Aprobar la captura o el relevo de identidad no basta para acreditar la continuidad completa del despacho.

La copia de stock, crédito y precios lleva fecha de actualización; al superar su vigencia se presenta como referencia, no como aprobación de la venta. El conductor tampoco presenta una captura POS como pago autorizado. Para el relevo de turno en bodega, la validación local de identidad debe probarse durante 24 horas antes de acreditar la autonomía completa.

